% ID: SEL-D0078-OA-E03
% Tanda S078.
\textbf{Parte I.} Suponemos un dado equilibrado y lanzamientos independientes. Ana gana si obtiene 6 en el primer lanzamiento:
\[
 P(\text{gana Ana})=\boxed{\frac16}.
\]
Para que gane Blas, Ana debe fallar y después Blas debe obtener 2 o 3:
\[
 P(\text{gana Blas})=\frac56\frac26=\boxed{\frac5{18}}.
\]
Para que no gane ninguno, deben fallar ambos:
\[
 P(\text{ninguno gana})=\frac56\frac46=\boxed{\frac59}.
\]
Son los tres resultados excluyentes de la partida, y sus probabilidades suman $1/6+5/18+5/9=1$. El juego termina tras un máximo de dos lanzamientos.

\textbf{Parte II.} La proporción muestral es $\hat p=450/1200=0{,}375$. Las frecuencias 450 y 750 permiten usar la aproximación normal. Para una confianza del $90\%$, queda $0{,}05$ en cada cola y $z_{0{,}95}\approx1{,}6449$. El intervalo es
\[
 0{,}375\pm1{,}6449\sqrt{\frac{0{,}375\cdot0{,}625}{1200}}
 \approx\boxed{[0{,}3520;\ 0{,}3980]}.
\]
